\documentclass[10pt,a4paper]{amsart}
\usepackage[margin=19mm]{geometry}
\usepackage{amsmath,amssymb,mathtools}
\usepackage[scr=rsfs,cal=euler]{mathalfa}
\usepackage{xfrac}
\usepackage[unicode,hidelinks,bookmarks=false]{hyperref}
\newcommand{\ie}{i.e.\ }
\newcommand{\cf}{cf.\ }
\newcommand{\resp}{respectively\ }
\newcommand{\Subsection}{\subsection}
\providecommand{\nobreakdash}{\nobreak\hbox{-}}
\setlength{\emergencystretch}{2em}
\multlinegap0pt
\usepackage{bm}
\usepackage{bbm}
\usepackage{dsfont}
\usepackage{xspace}
\usepackage{cancel}
\usepackage{mathtools}
\usepackage{amsthm}
\makeatletter
\@ifundefined{theorem}{\newtheorem{theorem}{Theorem}}{}
\@ifundefined{lemma}{\newtheorem{lemma}[theorem]{Lemma}}{}
\makeatother
\newcommand{\C}{\mathbb{C}}
\title[Average Winding Number for Determinantal Curves associated w]{Average Winding Number for Determinantal Curves associated with 2-Matrix Models in the Class AIII}
\author{Mathieu Yahiaoui, Mario Kieburg}
\date{Source: arXiv:2511.06669v2 (CC BY 4.0)}
\begin{document}
\maketitle

\noindent\emph{Source and licence.} Average Winding Number for Determinantal Curves associated with 2-Matrix Models in the Class AIII. Version arXiv:2511.06669v2 (CC BY 4.0), distributed under the Creative Commons Attribution 4.0 International licence, as recorded in the arXiv licence metadata for this exact version and shown on the abstract page https://arxiv.org/abs/2511.06669v2. Full rights notice: licenses/arxiv-2511.06669-CC-BY-4.0.txt.\par\smallskip

\noindent\textit{Editorial scope.} Counted unit: the Lemma whose source label is \texttt{lemma:existence} in the article (source0.tex lines 384--394, source label \texttt{lemma:existence}), delivered together with the article's own complete original proof of it (source0.tex lines 396--458, one \texttt{proof} environment of 63 lines). No mathematics is written, repaired or re-derived: statement and proof are the article's own text line for line. Required context retained uncounted: none. Every definition, display and label of those lines is retained unchanged. Omitted from this package and not redistributed: 1--383, 395--395, 459--1604. Nothing inside a retained excerpt is omitted or altered: every retained body is delivered exactly as the article's lines are written, so no omission receipt of any kind accompanies this package. Citations: the retained excerpts carry no citation command at all, so no bibliography entry of the article is transcribed and no reference is redistributed. Reference adaptation: none was needed and none was made; every reference inside the delivered unit resolves inside it, so no unresolved reference or citation is produced. Printed slips kept literal: no printed word, symbol, hypothesis or number of the counted statement or of its proof was altered. Layout and class: the article's own class is \texttt{sn-jnl} with packages including graphicx, multirow, amsmath, amssymb, amsfonts, amsthm, mathrsfs, appendix, textcomp, manyfoot, booktabs, algorithm, algorithmicx, algpseudocode; the wrapper is the lane's pinned \texttt{amsart} editorial wrapper at 10pt on A4, which restates the theorem-like environments the retained text needs and adds the article's own macro definitions that the retained text needs (\texttt{\textbackslash C}), with extra wrapper packages (bm, bbm, dsfont, xspace, cancel, mathtools). The delivered numbering is the unit's own. Assets: no retained span contains a figure environment, a graphics-inclusion command, a PDF-page inclusion, a table or a TikZ picture; no image file is copied into this package, the package declares no asset dependency and no image file is redistributed with it. Licence: version arXiv:2511.06669v2 of this article is distributed under the Creative Commons Attribution 4.0 International licence, as recorded in the arXiv licence metadata for this exact version and shown on the abstract page https://arxiv.org/abs/2511.06669v2. A full rights notice is delivered at licenses/arxiv-2511.06669-CC-BY-4.0.txt. Characters: every retained excerpt is pure ASCII, so the delivered engine, XeLaTeX with its default Latin Modern text font, reports no missing character.
\begin{lemma}\label{lemma:existence}
    Let $\omega$ be a P\'olya weight such that $\widetilde{\omega}$ is a P\'olya frequency function of order $2$, for any $j\in[0,2N-2]$ 
    \begin{equation}
        \int_{\C}|z|^{j}\omega(|z|^2)\mathrm{d}^2z<\infty.
    \end{equation}
    Let $m$ be a strictly positive integer, $(\alpha_{\ell})$ be $m$ complex numbers and $(\beta_{\ell})$ be $m$ pairwise distinct  non-zero complex numbers.
    Then, the following holds
    \begin{equation}
        \int_{\C^{N}}\left|\Delta_{N}(\mathbf{z})\right|^{2}\prod\limits_{j=1}^{N}\left(\omega(\left|z_{j}\right|^{2})\prod\limits_{\ell=1}^{m}\left|\frac{z_{j}-\alpha_{\ell}}{z_{j}-\beta_{\ell}}\right|\right)\mathrm{d}^{2}\mathbf{z}<+\infty.
    \end{equation}
\end{lemma}

\begin{proof}
    Making use of the triangular inequality, we have the following rough bound
    \begin{equation}
        \left|\Delta_{N}(\mathbf{z})\right|^{2}\leq \prod\limits_{j=1}^{N}(1+\left|z_{j}\right|)^{2(N-1)}.
    \end{equation}
    Let us fix $r>0$ small enough so that all closed discs $\overline{\mathrm{D}}(\beta_{j},r)\subset\C$ are pairwise disjoint and do not intersect with the origin. We then set
    \begin{equation}
        \Omega:=\left(\C\setminus \bigcup\limits_{j=1}^{m}\overline{\mathrm{D}}(\beta_{j},r)\right)^{N}.
    \end{equation}
    For $\mathbf{z}\in\Omega$, each coordinate satisfies $\left|z_{j}-\beta_{\ell}\right|\geq r$ for all $1\leq \ell\leq m$. Therefore, uniformly in $z_{j}$ where $j$ is fixed, we have
    \begin{equation}
        \prod\limits_{\ell=1}^{m}\left|\frac{z_{j}-\alpha_{\ell}}{z_{j}-\beta_{\ell}}\right|=\prod\limits_{\ell=1}^{m}\left|1+\frac{\beta_{\ell}-\alpha_{\ell}}{z_{j}-\beta_{\ell}}\right|\leq \prod\limits_{\ell=1}^{m}\left(1+\frac{|\beta_{\ell}-\alpha_{\ell}|}{r}\right).
    \end{equation}
    By assumption $\int_{\C}|z|^{j}\omega(|z|^2)\mathrm{d}^2z<\infty$ for any $j\in[0,2N-2]$, the integral 
    \begin{equation}
    \begin{split}
        &\int_{\Omega}\left|\Delta_{N}(\mathbf{z})\right|^{2}\prod\limits_{j=1}^{N}\left(\omega(\left|z_{j}\right|^{2})\prod\limits_{\ell=1}^{m}\left|\frac{z_{j}-\alpha_{\ell}}{z_{j}-\beta_{\ell}}\right|\right)\mathrm{d}^{2}\mathbf{z}
        \\
        \leq&\left(\int_{\C\setminus \bigcup\limits_{j=1}^{m}\overline{\mathrm{D}}(\beta_{j},r)}(1+\left|z\right|)^{2(N-1)}\prod\limits_{\ell=1}^{m}\left(1+\frac{|\beta_{\ell}-\alpha_{\ell}|}{r}\right)\omega(|z|^2)\mathrm{d}^2z\right)^N<\infty
    \end{split}
    \end{equation}
    exists on $\Omega$.


    On the complement $\Omega^{\complement}$, there exists an index pair $(j,\ell)$ such that $\left|z_{j}-\beta_{\ell}\right|< r$. We work inside the compact closed disc 
    \begin{equation}
        \overline{\mathrm{D}}(\beta_{\ell},r)=\left\{ z\in\mathbb{C}:\,\left| z-\beta_{\ell} \right|\leq r \right\}.
    \end{equation}
    Because $\omega$ is a P\'olya weight, it is positive, continuous and integrable on $\mathbb{R}_{+}$. In particular, any singularity must arise at a boundary point say $t_{0}\in\mathbb{R}_{+}$: there exists $\alpha\in[0,1[$ and $c_{1}>0$ such that, on a neighbourhood of $t_{0}$
    \begin{equation*}
        \omega(t)\leq c_{1}\left|t-t_{0}\right|^{-\alpha}
    \end{equation*}
    If $\left|\beta_{\ell}\right|^{2}\neq t_{0}$, we simply have $\alpha=0$ in what follows. Assume there exist an index $\ell$ such that $\left|\beta_{\ell}\right|^{2}= t_{0}$. Because $\beta_{\ell}$ is non-zero and the function $z\mapsto \left|z\right|^{2}$ has a non-vanishing gradient at $z=\beta_{\ell}$, for $r>0$ sufficiently small there exist constants $c_{2},c_{3}>0$ such that
    \begin{equation*}
        c_{2}\left|z-\beta_{\ell}\right|\leq \left|\left|z\right|^{2}-\left|\beta_{\ell}\right|^{2}\right| \leq c_{3}\left|z-\beta_{\ell}\right|.
    \end{equation*}
    Consequently, on $\overline{\mathrm{D}}(\beta_{\ell},r)$ there exists $c_{4}>0$ such that
    \begin{equation*}
        \omega(\left|z\right|^{2})\leq c_{4}\left|z-\beta_{\ell}\right|^{-\alpha}.
    \end{equation*}
    All factors in the integrand that remain away from their own singularities are bounded on this disc. In particular,
    \begin{itemize}
        \item all denominators $\left|z_{j}-\beta_{k}\right|$ with $k\neq \ell$ stay uniformly bounded away from zero since the discs are chosen disjoints,
        \item the Vandermonde term is bounded by
        \begin{equation*}
            \left|\Delta_{N}(\mathbf{z})\right|^{2}\leq \prod_{j=1}^{N}(1+\left|z_{j}\right|)^{2(N-1)}.
        \end{equation*}
    \end{itemize}
    It follows that, for $z_{j}\in \overline{\mathrm{D}}(\beta_{\ell},r)$, the factor
    \begin{equation*}
        (1+\left|z_{j}\right|)^{2(N-1)}\omega(\left|z_{j}\right|^{2})\left|z_{j}-\alpha_{\ell}\right|\prod\limits_{k\neq\ell}^{}\left|\frac{z_{j}-\alpha_{k}}{z_{j}-\beta_{k}}\right|
    \end{equation*}
    is bounded, up to a constant, by $\left|z_{j}-\beta_{\ell}\right|^{-\alpha}$. Hence, there exists $c_{5}>0$ such that the full integrand on $\overline{\mathrm{D}}(\beta_{\ell},r)$ is dominated by
    \begin{equation*}
        c_{5}\left|z_{j}-\beta_{\ell}\right|^{-(1+\alpha)}.
    \end{equation*}
    In polar coordinates around $\beta_{\ell}$, that is $z=\beta_{\ell}+\rho e^{i\theta}$, on this closed disc we have
    \begin{equation*}
        \int_{\overline{\mathrm{D}}(\beta_{\ell},r)}\left|z_{j}-\beta_{\ell}\right|^{-(1+\alpha)}\mathrm{d}^{2}z=\int_{0}^{2\pi}\mathrm{d}\theta\int_{0}^{r}\rho^{-(1+\alpha)}\rho\mathrm{d}\rho=2\pi\int_{0}^{r}\rho^{-\alpha}\mathrm{d}\rho <\infty,
    \end{equation*}
    since $\alpha<1$. Therefore, the integral over $\Omega^{\complement}$ is finite which allows us to conclude.
    
\end{proof}

\end{document}
